
% Page 54

novembre 1591. Le seigneur de Tanus avait d'ailleurs rejoint les assaillants. La ville d'Albi apprit avec grande consternation la nouvelle de cette capitulation. Les troupes royales s'étaient bien mises en marche, surprenant les réformés aux abords de Réalmont, mais Alban demeura aux mains des Protestants. Et cela dura jusqu'à l'Édit de Nantes en 1598.

Ainsi donc, en l'espace d'un quart de siècle, notre village changea une quinzaine de fois de mains... pillé tantôt par les uns, tantôt par les autres... Un temps où « les Français ne s'aimaient guère ».

1598, date de l'Édit de Nantes, met fin, au moins juridiquement, aux luttes fratricides. En octobre, Monseigneur DelBène peut faire son entrée dans sa ville épiscopale d'Albi, aux acclamations de ses diocésains. Ceux-ci l'attendent depuis dix ans.

On devine que dans notre coin de montagne, la misère est grande. Il est difficile de faire rentrer les impôts. Déjà en juin 1593, les consuls d'Alban font à qui de droit des déclarations en ce sens. En 1594, le sang coula. Le peuple s'opposait par la force aux saisies que les agents royaux voulaient mettre à exécution. Des gens d'un certain rang sont acculés à emprunter, ainsi un certain Manelphe, marchand demeurant à la Vialette et Cabannes habitant d'Alban. Le premier se fait prêter cent écus et le second vingt livres par un compatriote, Salvy Puech. Ce dernier possède pas mal de biens : une maison à deux étages (ce qui, pour l'époque, est un luxe), une grange, une étable, des prés, des bois, à la « Borie de Conte » au voisinage de Peyreblanque. Il a des terres aussi au Colombié avoisinant le pré de la Capélanyé (sous la nouvelle actuelle gendarmerie). Il possède aussi un jardin au foiral et un autre au Moulin du Vent. Ce riche propriétaire terrien a le goût des armes et fait partie de la compagnie du duc de Joyeuse.
